\begin{problem}[Distinguished-open transitivity]\label{prob:vi18-p19}
Starting with $D(f)\cong\Spec A_f$, prove that a distinguished open $D(g/1)$ inside it corresponds to $D(fg)$ inside $\Spec A$, including the sheaf identification.
\end{problem}

\ifdefined\FullProblemDossiers
\paragraph{Why this problem matters.}
This problem consolidates the iterated localization aspect of affine geometry and turns a structural statement into a reusable calculation.

\paragraph{Useful ideas before starting.}
Use $(A_f)_{g/1}\cong A_{fg}$.

\paragraph{Guided hints.}
\begin{enumerate}
\item Use $(A_f)_{g/1}\cong A_{fg}$.
\item Reduce the statement to a ring-theoretic property whenever possible.
\item Translate the result back into topology, sheaves, or local rings so that all affine-scheme layers are visible.
\end{enumerate}

\begin{solution}
Inside $\Spec A_f$, the distinguished open determined by $g/1$ has coordinate ring
\[
(A_f)_{g/1}.
\]
By associativity of localization,
\[
(A_f)_{g/1}\cong A_{fg}.
\]
Under the homeomorphism $\Spec A_f\cong D(f)$ from VI.18.P1, a prime $\mfp A_f$ lies in $D(g/1)$ iff $g/1\notin\mfp A_f$, equivalently $g\notin\mfp$.  Together with $f\notin\mfp$, this means $\mfp\in D(fg)$.  Thus
\[
D_{\Spec A_f}(g/1)\cong D_{\Spec A}(fg)\cong\Spec A_{fg},
\]
and the structure-sheaf identifications agree through the localization isomorphism.
\end{solution}

\paragraph{What happened mathematically.}
The calculation illustrates that global affine geometry is controlled by commutative algebra once the structure sheaf is kept in view.

\paragraph{Extensions.}
Vary the ring by products, quotients, localizations, arithmetic examples, or nilpotent thickenings and repeat the same dictionary.

\paragraph{Provenance.}
\textbf{New canonical transitivity variant of VI.18.P1; this is not a second migration of AG3 Problem 11.}
\else
\paragraph{Provenance.} \textbf{New canonical transitivity variant of VI.18.P1; this is not a second migration of AG3 Problem 11.}
\fi
